\documentclass[aspectratio=169,10pt,spanish]{beamer}

\input{../../../_preambulo_beamer.tex}
\input{../../../_comandos_beamer.tex}

\selectlanguage{spanish}

\title{MA1001B - Modelación Estadística para la Toma de Decisiones}
\subtitle{Lección 37: Dos medias independientes (t de dos muestras)}
\author{}
\institute{Tecnol\'ogico de Monterrey}
\date{}

\begin{document}

\begin{frame}[plain]
  \vspace{1.2cm}
  {\Large\bfseries MA1001B - Modelación Estadística para la Toma de Decisiones\par}
  \vspace{0.7cm}
  {\large Lección 37: Dos medias independientes (t de dos muestras)\par}
  \vspace{0.7cm}
  {\color{TecRojo}\rule{\textwidth}{1.2pt}}
  \vspace{0.8cm}

  Facultad MA1001B\\[0.3cm]
  Periodo\\[0.3cm]
  Tecnol\'ogico de Monterrey
\end{frame}

\begin{frame}{La pregunta de las medias de dos turnos}
  \begin{alertblock}{Pregunta guía}
    ¿Cómo se comparan los tiempos medios de recolección de dos turnos
    independientes cuando cada grupo tiene su propia $s$?
  \end{alertblock}

  \vspace{0.6em}
  Al final de esta lección, usted puede:
  \begin{itemize}
    \item mantener dos muestras independientes sin agrupar;
    \item calcular un estadístico $t$ de Welch y $gl$ de Satterthwaite;
    \item decidir $H_0:\mu_A=\mu_B$ a partir de un valor p bilateral;
    \item rechazar una lectura causal de un agrupamiento observacional.
  \end{itemize}

  \vfill
  \scriptsize Todos los tiempos de recolección y valores de experiencia usados hoy son sintéticos. No se usan datos reales de empresa.
\end{frame}

\begin{frame}{Dos muestras independientes, dos desviaciones estándar}
  \small
  Un almacén sintético muestrea $n_A=30$ recolecciones del turno A y
  $n_B=32$ recolecciones del turno B. Los asociados no están pareados, y las
  desviaciones estándar muestrales no tienen que ser iguales.

  \[
    t=\frac{\bar{x}_A-\bar{x}_B}{\sqrt{s_A^2/n_A+s_B^2/n_B}}.
  \]

  \begin{exampleblock}{Procedimiento declarado: Welch}
    $H_0:\mu_A=\mu_B$ frente a $H_1:\mu_A\neq\mu_B$, $\alpha=0.05$. Los
    grados de libertad son los de Satterthwaite, no $n_A+n_B-2$.
  \end{exampleblock}
\end{frame}

\begin{frame}[fragile]{Paso 1: Resúmenes de grupo}
  \begin{columns}[T]
    \begin{column}{0.42\textwidth}
      \small
      \begin{lstlisting}
mean_a = times_a.mean()
s_a = times_a.std(ddof=1)
diff = mean_a - mean_b
      \end{lstlisting}

      \begin{block}{Salida verificada}
        $\bar{x}_A=48.100885$, $s_A=4.660494$\\
        $\bar{x}_B=52.826941$, $s_B=6.325509$\\
        $\bar{x}_A-\bar{x}_B=-4.726057$
      \end{block}
    \end{column}
    \begin{column}{0.58\textwidth}
      \centering
      \includegraphics[width=\textwidth]{../figuras/t_dos_muestras_01_resumenes_grupos.png}
      \vspace{0.2em}
      \scriptsize Diagramas de caja de turnos independientes del Paso 1
    \end{column}
  \end{columns}
\end{frame}

\begin{frame}{La $t$ de Welch no agrupa las varianzas}
  \small
  \[
    EE=\sqrt{\frac{4.660494^2}{30}+\frac{6.325509^2}{32}}=1.405128,
  \]
  \[
    t=\frac{-4.726057}{1.405128}=-3.363436,\qquad gl=56.900423.
  \]
  El valor p bilateral es $0.001382$. Como $|t|>t^*=2.002541$ y $p<\alpha$,
  la prueba rechaza $H_0:\mu_A=\mu_B$.
\end{frame}

\begin{frame}[fragile]{Paso 2: Decisión de Welch de dos muestras}
  \begin{columns}[T]
    \begin{column}{0.40\textwidth}
      \small
      \begin{lstlisting}
se = sqrt(s_a**2/n_a + s_b**2/n_b)
t = (mean_a - mean_b) / se
p = 2 * t_dist.sf(abs(t), df)
      \end{lstlisting}

      \begin{block}{Salida verificada}
        $t=-3.363436$, $gl=56.900423$\\
        $p=0.001382$\\
        Decisión: rechazar $H_0$
      \end{block}
    \end{column}
    \begin{column}{0.60\textwidth}
      \centering
      \includegraphics[width=\textwidth]{../figuras/t_dos_muestras_02_prueba_welch.png}
      \vspace{0.2em}
      \scriptsize Colas $t$ de Welch del Paso 2
    \end{column}
  \end{columns}
\end{frame}

\begin{frame}{Paso 3: Los grupos observacionales no son un experimento}
  \begin{columns}[T]
    \begin{column}{0.42\textwidth}
      \small
      Experiencia media del turno A: $4.484559$ años.

      \vspace{0.3em}
      Experiencia media del turno B: $1.801265$ años.

      \vspace{0.4em}
      Los asociados no fueron asignados al azar a los turnos.

      \vspace{0.5em}
      \textbf{Límite:} un $p=0.001382$ de Welch significativo compara medias.
      No identifica el turno como una causa. La Lección 38 estudia pares
      emparejados antes-después.
    \end{column}
    \begin{column}{0.58\textwidth}
      \centering
      \includegraphics[width=\textwidth]{../figuras/t_dos_muestras_03_limite_observacional.png}
      \vspace{0.2em}
      \scriptsize Confounder de experiencia del Paso 3
    \end{column}
  \end{columns}
\end{frame}

\begin{frame}{Verificación de conceptos}
  \begin{enumerate}
    \item ¿Por qué estos dos turnos son muestras independientes y no pares?
    \item Escriba el error estándar de Welch usando $s_A$, $s_B$, $n_A$ y
      $n_B$.
    \item Si $p=0.001382$ y $\alpha=0.05$, ¿cuál es la decisión para
      $H_0:\mu_A=\mu_B$?
    \item ¿Por qué una diferencia de experiencia media bloquea una afirmación
      causal?
  \end{enumerate}

  \vspace{0.6em}
  \begin{block}{Conexión con la Lección 38}
    La sesión puede continuar directamente con una prueba $t$ pareada, donde
    las mismas unidades se miden antes y después de un cambio.
  \end{block}

  \vfill
  \scriptsize Fuente: OpenStax, \textit{Introductory Business Statistics 2e},
  Capítulo 10, Sección 10.1. CC BY-NC-SA.
\end{frame}

\appendix



\end{document}
